\subsection*{从完整预期开始构造历史}

先为一条历史写下初态和预期，再实现执行器。章首的单次外部购买，预期包含一个 $S$ 绑定、P1/C1/T1、平台 X1、余额 7000、原 O1/EV1、消费者 F501 和固定 dispatch。各表中的账户、币种、金额、引用、原输入与结果都要一致。仅检查订单数为一，会漏掉把 S1 错配到 S2 的订单这种错误。

原结果也要独立指定。例如 \code{payment\_result} 中的关联与金额为：
\begin{BookCode}
purchase_id=P1, charge_id=C1, order_id=O1
provider_charge_id=X1, amount_minor=3000
\end{BookCode}
EV1 载荷引用同一组编号，\code{occurred\_at} 是决定首次形成时保存的原值。F1 的实际编号从消费插入取得，为 501，业务回报使用这个实际编号。

然后把故障点命名。例如 $A_J$ 提交后本地 A 以前退出，平台 E 提交后 $R_J$ 以前丢响应，$R_J$ 提交后本地 R 以前退出，EV1 发布后标记以前退出，以及 Q 提交后 ACK 以前退出。每次退出都只改变允许丢失的进程状态，检查者从继续保存的事实推导恢复动作。

\CFifteenFigure{ch15-06-faults}{故障位于具体提交之间，恢复从已保存的事实接续尚未完成的责任}{fig:ch15-faults}

一次端到端历史可以依次经过多个故障，检查者也保留每一步的完整状态。平台成功、本地未决是合法组合，本地成功却没有可核对平台证据则违反结果真实性。检查器按各时刻已经提交的事实判断，允许恢复过程经过这些跨域过渡。

\subsection*{允许迁移如何成为检查器}

为操作建立顺序规格，规定首次受理、同参恢复、异参冲突、成功证据合流和最终拒绝分别能改变什么。检查执行历史时，把每项原子决定放到与实际提交相容的位置。两个重叠的首次请求可以由任意一个先提交，后来的来件根据胜者的输入与结果处理。响应丢失只删除调用方的获知事件，已提交事实继续存在。

例如，同 $S$ 的 A 输入事务先开始却回滚，B 输入随后取得首次已提交绑定，A 恢复时得到冲突。这条历史符合首次提交绑定方案。检查器保护 B 的绑定不可覆盖，并要求 A 的恢复没有形成另一项效果。若业务要求最先到达的输入永远胜出，就要增加提前受理的权威决定和相应恢复责任。

并发测试据此控制关键交错。两个不同 $S$、相同 3000 输入都应保留各自购买。两个相同 $S$、不同输入只允许一份正式绑定。工作者 W2 取得新代次以后，W1 的本地完成更新应返回零行或明确旧资格拒绝，平台调用则继续用原 \code{pay-C1}。预期集合先给出允许的胜者与结果，执行者依据这组预期判断本次历史。

\subsection*{有限模型用来找哪类反例}

退款提供了一个小而完整的模型。原收款额取 3 个抽象金额单位，设置 RF1、RF2、RF3，每项金额可取 1 或 2。每项身份处于空、未决、成功、最终拒绝之一。空身份受理时预占金额，同参重试和未知观察保持状态，未决得到成功证据后转为成功，取得足够拒绝证据后释放占用。

检查器遍历这些允许迁移，始终检查成功金额与未决占用之和不超过 3，并检查每个 RF 的金额一旦绑定就保持。这里每一金额单位对应 1000 分，容量 3 对应原收款 3000 分。枚举的对象是上述三个身份、两种金额和四种阶段组成的有限状态图。

\CFifteenFigure{ch15-07-model}{退款有限模型固定金额与身份，成功转移占用，足够拒绝释放占用}{fig:ch15-model}

把受理条件故意改为只统计成功退款，就会出现很短的反例：RF1 预占 2，尚未成功；RF2 又预占 2。两次首次判断各看见成功总额为 0，都通过，未决责任合计成为 4。随后两项都成功，会超过原收款 3。反例把缺口定位在首次受理的额度仲裁，查询次数增加后，这个错误守卫也会作出相同决定。

纸面归纳说明任意一项允许迁移怎样保持关系，有限枚举帮助发现遗漏迁移、错误守卫和状态组合。数据库实验再检查所选 SQL、隔离级别与条件更新是否落实了这些守卫。三者使用一致的身份与状态定义，就可以沿同一分支核对证明、模型和实现。

\subsection*{恢复演练还要破坏一行已有数据}

从平台 X1、O1、EV1、F501 都已经形成的成功状态开始，恢复购买库的旧逻辑备份，让本地 P/C/O/EV 消失，其他保存域保持。恢复者根据独立 J 重建原关系，再重投原 EV1。消费者应找回 F501，平台余额继续为 7000。

这条历史检查空位置重建。再给本地库留下一行损坏的 O1，例如主键叫 O1，正文却引用 \code{P-other}。恢复者先发现 T1 缺失并准备补写，随后核对 O1 发现归属冲突。若每张表分别提交，T1 会留下，本批物化边界也可能向前推进，形成半批恢复。

恢复事务因此把字段比较、补写、允许的状态推进和物化边界共同提交。已有行等于原 $A_J$ 的未决投影，可以按精确条件推进；已经等于 $R_J$ 的目标投影，可以确认已有内容；其他内容进入冲突。完整原结果、事件版本、原 \code{occurred\_at} 和全部 payload 都参与比较。

\begin{BookCode}
取得已核验的同一权威快照与目标边界
BEGIN 本地物化事务
锁定本分区恢复资格，核验目标代次
逐项读取原 P/C/T/O/EV 和其关联
已有行：
    逐字段比较原归属、输入、平台关联、原结果与事件
    等于允许的原未决投影：按完整旧值条件推进
        影响行数必须为 1，否则抛错
    等于本次目标投影：保留已有内容
    其他情况：抛出恢复冲突
缺失行：
    插入权威决定中的原实际编号与完整内容
    唯一约束冲突：抛错，不覆盖胜出行
共同保存本批物化边界
COMMIT
提交成功后再发布完成
任一步抛错：整批回滚，在事务外保存核查记录并保持受控入口
提交未知：重读原编号及边界，按同一规则恢复
\end{BookCode}

在 PostgreSQL 的 READ COMMITTED 下，\code{FOR UPDATE} 等待竞争者结束，再锁定并返回符合条件的更新后行，每条命令取得自己的可见快照~\cite{pg}。物化者还要固定本批权威输入。读取 J 的一半结果以后允许另一事务继续分配新 O/EV，再把另一半拼进来，就会混合两个边界。恢复期间冻结相应写入，或者从 J 的同一完整快照读取，给出了可解释的目标图像。物化事务取得本地资格锁后还会核验代次，目标来源和本地写资格各自有检查位置。

\CFifteenFigure{ch15-08-rollback}{T1 的补写与损坏 O1 的核对同属一个事务，冲突使数据与物化边界一起回滚}{fig:ch15-rollback}

检查这条反例时，保留冲突前后所有相关表的完整内容，比较原绑定、结果、事件以及游标。还要把已有订单改为 O2，让它与待恢复的 O1 争用同一个 C1，或争用同一平台范围内的 X1。两个对象编号不同，业务唯一索引应使后插入者失败，撤销本批先补入的 T、EV 和台账，游标停在原位置。另取两个合法购买，分别使用自己的 P/C/T 和不同平台收款，确认这些约束允许它们共同保存。再制造一次本应更新一行却实际为零行的分支，确认之前已经准备的 P 进度更新也回滚。异常记录说明控制流进入了失败出口，前后完整快照则说明这一出口撤销了哪些准备写入。

\subsection*{证据范围与工程结果怎样对应}

真实数据库事务、受控平台接口、有限状态枚举和复制产品实验分别登记。数据库在具体版本与配置下执行行锁、约束和回滚，检验相应实现路径。平台夹具按给定规则保存键和资金账本，检验这组规则下的协议历史。消息代理与共识组件的实际部署，则由相应版本、配置和原始日志连接到其故障模型。

整链路采用受控投递重复 EV1 时，检查对象是 Q 的事务与恢复。消息代理的持久化、确认和重投规则沿第九、十章的具体链路核对：发布者得到代理确认后才标记本地发布进度，消费者提交 Q 后才 ACK，代理与接收者之间的重投继续携带 EV1。第十一章的复制实验把三份独立进程的日志、提交、快照与切换接到状态机应用边界。本章集成路径中的 J 则由一个 PostgreSQL 事务域提供权威仲裁。证据索引分别记录这些对象，修改某一迁移时就能找到相应运行。
\FloatBarrier
